\section{АНАЛИТИЧЕСКИЙ ОБЗОР}

В разделе рассматриваются эволюция систем хранения аналитических данных, открытые
табличные форматы как ключевой компонент архитектуры lakehouse, колоночные файловые
форматы хранения, архитектура системы Apache Paimon и применяемые форматами
алгоритмы кодирования, сжатия и пропуска данных. По итогам обзора формулируется
постановка задачи.

\subsection{Эволюция систем хранения аналитических данных}

Исторически аналитическая обработка данных опиралась на хранилища данных~\cite{inmon2005datawarehouse, kimball2013toolkit} (data
warehouse) --- специализированные реляционные системы, в которые данные загружаются
по расписанию через процессы ETL, приводятся к единой схеме «звезда» или
«снежинка» и затем обслуживают запросы бизнес-аналитики. Хранилище данных
обеспечивает строгую типизацию, управление схемой, транзакционность и
предсказуемую производительность запросов. Платой за это являются высокая стоимость
хранения (вычислительные узлы и хранилище связаны и масштабируются совместно),
жёсткость схемы (изменение требует перепроектирования ETL) и неспособность работать
со слабоструктурированными данными --- журналами, событиями, телеметрией.

С появлением распределённой файловой системы Hadoop (HDFS) сформировалась
альтернативная архитектура --- озеро данных (data lake). Данные складываются в HDFS
или объектное хранилище «как есть», в исходных форматах, а схема применяется во
время чтения (schema-on-read). Слой совместимости с SQL обеспечивался Hive Metastore~\cite{thusoo2009hive}
--- сервисом каталога, хранящим описания таблиц (расположение файлов, схему,
разбиение на разделы) в реляционной базе, --- и движками SQL поверх него: Apache
Hive~\cite{thusoo2009hive}, затем Apache Impala~\cite{kornacker2015impala}, Presto, Spark SQL. Озеро данных дёшево и масштабируемо,
позволяет хранить любые данные, но не даёт транзакционных гарантий: одновременная
запись и чтение, частичные обновления, атомарная замена набора файлов не
поддерживаются на уровне файловой системы. Hive-таблица, по сути, --- это директория
с файлами; согласованность при изменениях обеспечивается только дисциплиной на
стороне приложений.

Ограничения обоих подходов привели к появлению архитектуры lakehouse. Идея состоит
в том, чтобы сохранить дешёвое масштабируемое объектное хранилище озера данных, но
добавить поверх него слой метаданных, который превращает набор файлов в полноценную
таблицу с ACID-транзакциями, эволюцией схемы, моментальными снимками (snapshot
isolation), управлением версиями и параллельным доступом. Этот слой и называется
\textit{открытым табличным форматом}. Вычислительные движки (Spark, Flink, Trino и
другие) работают с таблицей через спецификацию табличного формата, а не напрямую с
файлами; разделение хранения и вычислений сохраняется, поскольку и данные, и
метаданные лежат в объектном хранилище. Современные реализации --- Apache Iceberg~\cite{iceberg-spec},
Delta Lake~\cite{armbrust2020delta}, Apache Hudi~\cite{hudi-docs} и Apache Paimon~\cite{paimon-docs} --- различаются деталями устройства
метаданных и набором поддерживаемых сценариев, но разделяют общую модель.

Особое место среди них занимает Apache Paimon. В отличие от Iceberg и Delta Lake,
ориентированных в первую очередь на таблицы только для добавления (append-only) и
пакетную загрузку, Paimon изначально проектировался под потоковую обработку и
таблицы с первичным ключом: в его основе лежит LSM-дерево, что делает эффективными
сценарии захвата изменений (Change Data Capture) и построения near-real-time-витрин.
Именно поэтому Apache Paimon выбран целевой системой настоящей работы.
